\section{The \texttt{Mirror\_toroid\_pothole} McXtrace Component}
Toroidal shape mirror (in XZ)

\subsection*{Identification}
\begin{itemize}
  \item \textbf{Author:} Erik B Knudsen
  \item \textbf{Origin:} DTU Physics
  \item \textbf{Date:} Jul 2016
\end{itemize}

\subsection*{Description}
This is an implementation of a toroidal mirror which may be curved in two dimensions. To avoid solving quartic equations, the intersection is computed as a combination of two intersections. First, the ray is intersected with a cylinder to catch (almost) the small radius curvature. Secondly, the ray is the intersected with an ellipsoid, with the curvatures matching that of the torus.

The first incarnation (Mirror\_toroid.comp) the mirror curves outwards (a bump), but this incarnation (Mirror\_toroid\_pothole) curves inwards (a pothole).

Example: Mirror\_toroid\_pothole( radius=0.1, radius\_o=1000, xwidth=5e-2, zdepth=2e-1,R0=1, coating="")

\subsection*{Input parameters}
Parameters in \textbf{boldface} are required; the others are optional.

\begin{longtable}{p{0.22\textwidth}p{0.12\textwidth}p{0.46\textwidth}p{0.14\textwidth}}
\toprule
\textbf{Name} & \textbf{Unit} & \textbf{Description} & \textbf{Default} \\
\midrule
\endhead
zdepth & m & Length of mirror. & 0.1 \\
xwidth & m & Width of mirror. & 0.01 \\
\textbf{radius} & m & Curvature radius &  \\
\textbf{radius\_o} & m & Curvature radius, outwards &  \\
R0 & 1 & Reflectivity of mirror. & 0 \\
coating & str & Datafile containing either mirror material constants or reflectivity  numbers. & "" \\
\bottomrule
\end{longtable}

\subsection*{Links}
\begin{itemize}
  \item Component source code found in file \texttt{Mirror\_toroid\_pothole.comp}.
\end{itemize}
\IfFileExists{contrib/Mirror_toroid_pothole_static.tex}{\input{contrib/Mirror_toroid_pothole_static.tex}}{}